\section{Relative Proj as a functor}
\label{section-relative-proj}

\noindent
We place ourselves in Situation \ref{situation-relative-proj}.
So $S$ is a scheme and $\mathcal{A} = \bigoplus_{d \geq 0} \mathcal{A}_d$
is a quasi-coherent graded $\mathcal{O}_S$-algebra.
In this section we relativize the construction of
$\text{Proj}$ by constructing a functor which the relative
homogeneous spectrum will represent.
As a result we will construct a morphism of schemes
$$
\underline{\text{Proj}}_S(\mathcal{A}) \longrightarrow S
$$
which above affine opens of $S$ will look like the homogeneous spectrum
of a graded ring. The discussion will be modeled after our
discussion of the relative spectrum in Section \ref{section-spec}.
The easier method using glueing schemes of the form
$\text{Proj}(A)$, $A = \Gamma(U, \mathcal{A})$, $U \subset S$
affine open, is explained in Section \ref{section-relative-proj-via-glueing},
and the result in this section will be shown to be isomorphic to that one.

\medskip\noindent
Fix for the moment an integer $d \geq 1$.
We denote $\mathcal{A}^{(d)} = \bigoplus_{n \geq 0} \mathcal{A}_{nd}$
similarly to the notation in Algebra, Section \ref{algebra-section-graded}.
Let $T$ be a scheme.
Let us consider {\it quadruples $(d, f : T \to S, \mathcal{L}, \psi)$
over $T$} where
\begin{enumerate}
\item $d$ is the integer we fixed above,
\item $f : T \to S$ is a morphism of schemes,
\item $\mathcal{L}$ is an invertible $\mathcal{O}_T$-module, and
\item
$\psi : f^*\mathcal{A}^{(d)} \to \bigoplus_{n \geq 0}\mathcal{L}^{\otimes n}$
is a homomorphism of graded $\mathcal{O}_T$-algebras
such that $f^*\mathcal{A}_d \to \mathcal{L}$ is surjective.
\end{enumerate}
Given a morphism $h : T' \to T$ and a quadruple
$(d, f, \mathcal{L}, \psi)$ over $T$ we can pull it back to the
quadruple $(d, f \circ h, h^*\mathcal{L}, h^*\psi)$ over $T'$.
Given two quadruples $(d, f, \mathcal{L}, \psi)$ and
$(d, f', \mathcal{L}', \psi')$ over $T$ with the same integer $d$
we say they are {\it strictly equivalent} if $f = f'$ and there exists
an isomorphism $\beta : \mathcal{L} \to \mathcal{L}'$
such that $\beta \circ \psi = \psi'$ as graded $\mathcal{O}_T$-algebra maps
$f^*\mathcal{A}^{(d)} \to \bigoplus_{n \geq 0} (\mathcal{L}')^{\otimes n}$.

\medskip\noindent
For each integer $d \geq 1$ we define
\begin{eqnarray*}
F_d : \Sch^{opp} & \longrightarrow & \textit{Sets}, \\
T & \longmapsto &
\{\text{strict equivalence classes of }
(d, f : T \to S, \mathcal{L}, \psi)
\text{ as above}\}
\end{eqnarray*}
with pullbacks as defined above.

\begin{lemma}
\label{lemma-proj-base-change}
In Situation \ref{situation-relative-proj}. Let $d \geq 1$.
Let $F_d$ be the functor
associated to $(S, \mathcal{A})$ above.
Let $g : S' \to S$ be a morphism of schemes.
Set $\mathcal{A}' = g^*\mathcal{A}$. Let $F_d'$ be the
functor associated to $(S', \mathcal{A}')$ above.
Then there is a canonical isomorphism
$$
F'_d \cong h_{S'} \times_{h_S} F_d
$$
of functors.
\end{lemma}

\begin{proof}
A quadruple
$(d, f' : T \to S', \mathcal{L}',
\psi' : (f')^*(\mathcal{A}')^{(d)} \to
\bigoplus_{n \geq 0} (\mathcal{L}')^{\otimes n})$
is the same as a quadruple
$(d, f, \mathcal{L},
\psi : f^*\mathcal{A}^{(d)} \to
\bigoplus_{n \geq 0} \mathcal{L}^{\otimes n})$
together with a factorization of $f$ as $f = g \circ f'$. Namely,
the correspondence is $f = g \circ f'$, $\mathcal{L} = \mathcal{L}'$
and $\psi = \psi'$ via the identifications
$(f')^*(\mathcal{A}')^{(d)} = (f')^*g^*(\mathcal{A}^{(d)}) =
f^*\mathcal{A}^{(d)}$. Hence the lemma.
\end{proof}

\begin{lemma}
\label{lemma-relative-proj-affine}
In Situation \ref{situation-relative-proj}. Let $F_d$ be the functor
associated to $(d, S, \mathcal{A})$ above.
If $S$ is affine, then $F_d$ is representable by the open subscheme
$U_d$ (\ref{equation-Ud})
of the scheme $\text{Proj}(\Gamma(S, \mathcal{A}))$.
\end{lemma}

\begin{proof}
Write $S = \Spec(R)$ and $A = \Gamma(S, \mathcal{A})$.
Then $A$ is a graded $R$-algebra and $\mathcal{A} = \widetilde A$.
To prove the lemma we have to identify the functor $F_d$
with the functor $F_d^{triples}$ of triples defined in Section
\ref{section-morphisms-proj}.

\medskip\noindent
Let $(d, f : T \to S, \mathcal{L}, \psi)$ be a quadruple.
We may think of $\psi$ as a $\mathcal{O}_S$-module map
$\mathcal{A}^{(d)} \to \bigoplus_{n \geq 0} f_*\mathcal{L}^{\otimes n}$.
Since $\mathcal{A}^{(d)}$ is quasi-coherent this is the same
thing as an $R$-linear homomorphism of graded rings
$A^{(d)} \to \Gamma(S, \bigoplus_{n \geq 0} f_*\mathcal{L}^{\otimes n})$.
Clearly, $\Gamma(S, \bigoplus_{n \geq 0} f_*\mathcal{L}^{\otimes n}) =
\Gamma_*(T, \mathcal{L})$. Thus we may associate to
the quadruple the triple $(d, \mathcal{L}, \psi)$.

\medskip\noindent
Conversely, let $(d, \mathcal{L}, \psi)$ be a triple.
The composition $R \to A_0 \to \Gamma(T, \mathcal{O}_T)$
determines a morphism $f : T \to S = \Spec(R)$, see
Schemes, Lemma \ref{schemes-lemma-morphism-into-affine}.
With this choice of $f$ the map
$A^{(d)} \to \Gamma(S, \bigoplus_{n \geq 0} f_*\mathcal{L}^{\otimes n})$
is $R$-linear, and hence corresponds to a $\psi$ which we
can use for a quadruple $(d, f : T \to S, \mathcal{L}, \psi)$.
We omit the verification that this establishes an isomorphism
of functors $F_d = F_d^{triples}$.
\end{proof}

\begin{lemma}
\label{lemma-relative-proj-d}
In Situation \ref{situation-relative-proj}.
The functor $F_d$ is representable by a scheme.
\end{lemma}

\begin{proof}
We are going to use Schemes, Lemma \ref{schemes-lemma-glue-functors}.

\medskip\noindent
First we check that $F_d$ satisfies the sheaf property for the
Zariski topology. Namely, suppose that $T$ is a scheme,
that $T = \bigcup_{i \in I} U_i$ is an open covering,
and that $(d, f_i, \mathcal{L}_i, \psi_i) \in F_d(U_i)$ such that
$(d, f_i, \mathcal{L}_i, \psi_i)|_{U_i \cap U_j}$ and
$(d, f_j, \mathcal{L}_j, \psi_j)|_{U_i \cap U_j}$ are strictly
equivalent. This implies that the morphisms $f_i : U_i \to S$
glue to a morphism of schemes $f : T \to S$ such that
$f|_{I_i} = f_i$, see Schemes, Section \ref{schemes-section-glueing-schemes}.
Thus $f_i^*\mathcal{A}^{(d)} = f^*\mathcal{A}^{(d)}|_{U_i}$.
It also implies there exist isomorphisms
$\beta_{ij} : \mathcal{L}_i|_{U_i \cap U_j} \to \mathcal{L}_j|_{U_i \cap U_j}$
such that $\beta_{ij} \circ \psi_i = \psi_j$ on $U_i \cap U_j$.
Note that the isomorphisms $\beta_{ij}$ are uniquely determined
by this requirement because the maps $f_i^*\mathcal{A}_d \to \mathcal{L}_i$
are surjective. In particular we see that
$\beta_{jk} \circ \beta_{ij} = \beta_{ik}$ on $U_i \cap U_j \cap U_k$.
Hence by Sheaves,
Section \ref{sheaves-section-glueing-sheaves} the invertible sheaves
$\mathcal{L}_i$ glue to an invertible $\mathcal{O}_T$-module
$\mathcal{L}$ and the morphisms $\psi_i$ glue to
morphism of $\mathcal{O}_T$-algebras
$\psi : f^*\mathcal{A}^{(d)} \to \bigoplus_{n \geq 0} \mathcal{L}^{\otimes n}$.
This proves that $F_d$ satisfies the sheaf condition with respect to
the Zariski topology.

\medskip\noindent
Let $S = \bigcup_{i \in I} U_i$ be an affine open covering.
Let $F_{d, i} \subset F_d$ be the subfunctor consisting of
those pairs $(f : T \to S, \varphi)$ such that
$f(T) \subset U_i$.

\medskip\noindent
We have to show each $F_{d, i}$ is representable.
This is the case because $F_{d, i}$ is identified with
the functor associated to $U_i$ equipped with
the quasi-coherent graded $\mathcal{O}_{U_i}$-algebra
$\mathcal{A}|_{U_i}$ by Lemma \ref{lemma-proj-base-change}.
Thus the result follows from Lemma \ref{lemma-relative-proj-affine}.

\medskip\noindent
Next we show that $F_{d, i} \subset F_d$ is representable by open immersions.
Let $(f : T \to S, \varphi) \in F_d(T)$. Consider $V_i = f^{-1}(U_i)$.
It follows from the definition of $F_{d, i}$ that given $a : T' \to T$
we have $a^*(f, \varphi) \in F_{d, i}(T')$ if and only if $a(T') \subset V_i$.
This is what we were required to show.

\medskip\noindent
Finally, we have to show that the collection $(F_{d, i})_{i \in I}$
covers $F_d$. Let $(f : T \to S, \varphi) \in F_d(T)$.
Consider $V_i = f^{-1}(U_i)$. Since $S = \bigcup_{i \in I} U_i$
is an open covering of $S$ we see that $T = \bigcup_{i \in I} V_i$
is an open covering of $T$. Moreover $(f, \varphi)|_{V_i} \in F_{d, i}(V_i)$.
This finishes the proof of the lemma.
\end{proof}

\noindent
At this point we can redo the material at the end of
Section \ref{section-morphisms-proj} in the current
relative setting and define a functor
which is representable by
$\underline{\text{Proj}}_S(\mathcal{A})$. To do this we introduce the
notion of equivalence between two
quadruples $(d, f : T \to S, \mathcal{L}, \psi)$ and
$(d', f' : T \to S, \mathcal{L}', \psi')$ with possibly different
values of the integers $d, d'$. Namely, we say these
are {\it equivalent} if $f = f'$, and there exists an
isomorphism $\beta : \mathcal{L}^{\otimes d'} \to (\mathcal{L}')^{\otimes d}$
such that
$\beta \circ \psi|_{f^*\mathcal{A}^{(dd')}} = \psi'|_{f^*\mathcal{A}^{(dd')}}$.
The following lemma implies that this defines an equivalence relation.
(This is not a complete triviality.)

\begin{lemma}
\label{lemma-equivalent-relative}
In Situation \ref{situation-relative-proj}.
Let $T$ be a scheme.
Let $(d, f, \mathcal{L}, \psi)$, $(d', f', \mathcal{L}', \psi')$
be two quadruples over $T$. The following are equivalent:
\begin{enumerate}
\item Let $m = \text{lcm}(d, d')$. Write $m = ad = a'd'$.
We have $f = f'$ and there exists
an isomorphism
$\beta : \mathcal{L}^{\otimes a} \to (\mathcal{L}')^{\otimes a'}$
with the property that $\beta \circ \psi|_{f^*\mathcal{A}^{(m)}}$
and $\psi'|_{f^*\mathcal{A}^{(m)}}$ agree
as graded ring maps
$f^*\mathcal{A}^{(m)} \to \bigoplus_{n \geq 0} (\mathcal{L}')^{\otimes mn}$.
\item The quadruples $(d, f, \mathcal{L}, \psi)$ and
$(d', f', \mathcal{L}', \psi')$ are equivalent.
\item We have $f = f'$ and
for some positive integer $m = ad = a'd'$ there exists an isomorphism
$\beta : \mathcal{L}^{\otimes a} \to (\mathcal{L}')^{\otimes a'}$
with the property that $\beta \circ \psi|_{f^*\mathcal{A}^{(m)}}$
and $\psi'|_{f^*\mathcal{A}^{(m)}}$ agree
as graded ring maps
$f^*\mathcal{A}^{(m)} \to \bigoplus_{n \geq 0} (\mathcal{L}')^{\otimes mn}$.
\end{enumerate}
\end{lemma}

\begin{proof}
Clearly (1) implies (2) and (2) implies (3) by restricting to
more divisible degrees and powers of invertible sheaves.
Assume (3) for some integer $m = ad = a'd'$. Let
$m_0 = \text{lcm}(d, d')$ and write it as $m_0 = a_0d = a'_0d'$.
We are given an isomorphism
$\beta : \mathcal{L}^{\otimes a} \to (\mathcal{L}')^{\otimes a'}$
with the property described in (3). We want to find an isomorphism
$\beta_0 : \mathcal{L}^{\otimes a_0} \to (\mathcal{L}')^{\otimes a'_0}$
having that property as well.
Since by assumption the maps $\psi : f^*\mathcal{A}_d \to \mathcal{L}$
and $\psi' : (f')^*\mathcal{A}_{d'} \to \mathcal{L}'$ are surjective the
same is true for the maps
$\psi : f^*\mathcal{A}_{m_0} \to \mathcal{L}^{\otimes a_0}$
and $\psi' : (f')^*\mathcal{A}_{m_0} \to (\mathcal{L}')^{\otimes a_0}$.
Hence if $\beta_0$ exists it is uniquely determined by the
condition that $\beta_0 \circ \psi = \psi'$. This means that
we may work locally on $T$. Hence we may assume that
$f = f' : T \to S$ maps into an affine open, in other words
we may assume that $S$ is affine. In this case the result follows
from the corresponding result for triples (see Lemma \ref{lemma-equivalent})
and the fact that triples and quadruples correspond in the
affine base case (see proof of Lemma \ref{lemma-relative-proj-affine}).
\end{proof}

\noindent
Suppose $d' = ad$. Consider the transformation of functors $F_d \to F_{d'}$
which assigns to the quadruple $(d, f, \mathcal{L}, \psi)$ over
$T$ the quadruple
$(d', f, \mathcal{L}^{\otimes a}, \psi|_{f^*\mathcal{A}^{(d')}})$.
One of the implications of Lemma \ref{lemma-equivalent-relative} is that the
transformation $F_d \to F_{d'}$ is injective!
For a quasi-compact scheme $T$ we define
$$
F(T) = \bigcup\nolimits_{d \in \mathbf{N}} F_d(T)
$$
with transition maps as explained above. This clearly defines a
contravariant functor on the category of quasi-compact schemes
with values in sets. For a general scheme
$T$ we define
$$
F(T)
=
\lim_{V \subset T\text{ quasi-compact open}} F(V).
$$
In other words, an element $\xi$ of $F(T)$ corresponds to a compatible system
of choices of elements $\xi_V \in F(V)$ where $V$ ranges over the
quasi-compact opens of $T$.
We omit the definition of the pullback map $F(T) \to F(T')$
for a morphism $T' \to T$ of schemes.
Thus we have defined our functor
\begin{equation}
\label{equation-proj}
F : \Sch^{opp} \longrightarrow \textit{Sets}
\end{equation}

\begin{lemma}
\label{lemma-relative-proj}
In Situation \ref{situation-relative-proj}.
The functor $F$ above is representable by a scheme.
\end{lemma}

\begin{proof}
Let $U_d \to S$ be the scheme representing the functor $F_d$
defined above. Let $\mathcal{L}_d$,
$\psi^d : \pi_d^*\mathcal{A}^{(d)} \to
\bigoplus_{n \geq 0} \mathcal{L}_d^{\otimes n}$ be the universal object.
If $d | d'$, then we may consider the quadruple
$(d', \pi_d, \mathcal{L}_d^{\otimes d'/d}, \psi^d|_{\mathcal{A}^{(d')}})$
which determines a canonical morphism $U_d \to U_{d'}$ over $S$.
By construction this morphism corresponds to the transformation
of functors $F_d \to F_{d'}$ defined above.

\medskip\noindent
For every affine open $\Spec(R) = V \subset S$
setting $A = \Gamma(V, \mathcal{A})$ we have a canonical
identification of the base change $U_{d, V}$ with the corresponding open
subscheme of $\text{Proj}(A)$, see Lemma \ref{lemma-relative-proj-affine}.
Moreover, the morphisms $U_{d, V} \to U_{d', V}$ constructed above
correspond to the inclusions of opens in $\text{Proj}(A)$.
Thus we conclude that $U_d \to U_{d'}$ is an open immersion.

\medskip\noindent
This allows us to construct $X$
by glueing the schemes $U_d$ along the open immersions $U_d \to U_{d'}$.
Technically, it is convenient to choose a sequence
$d_1 | d_2 | d_3 | \ldots$ such that every positive integer
divides one of the $d_i$ and to simply take
$X = \bigcup U_{d_i}$ using the open immersions above.
It is then a simple matter to prove that $X$ represents the
functor $F$.
\end{proof}

\begin{lemma}
\label{lemma-glueing-gives-functor-proj}
In Situation \ref{situation-relative-proj}.
The scheme $\pi : \underline{\text{Proj}}_S(\mathcal{A}) \to S$
constructed in Lemma \ref{lemma-glue-relative-proj}
and the scheme representing the functor $F$
are canonically isomorphic as schemes over $S$.
\end{lemma}

\begin{proof}
Let $X$ be the scheme representing the functor $F$.
Note that $X$ is a scheme over $S$ since the functor $F$
comes equipped with a natural transformation $F \to h_S$.
Write $Y = \underline{\text{Proj}}_S(\mathcal{A})$.
We have to show that $X \cong Y$ as $S$-schemes.
We give two arguments.

\medskip\noindent
The first argument uses the construction of $X$ as the union
of the schemes $U_d$ representing $F_d$ in the
proof of Lemma \ref{lemma-relative-proj}.
Over each affine open of $S$ we can identify $X$ with the homogeneous spectrum
of the sections of $\mathcal{A}$ over that open, since this was
true for the opens $U_d$. Moreover, these identifications
are compatible with further restrictions to smaller affine opens.
On the other hand, $Y$ was constructed by glueing these
homogeneous spectra.
Hence we can glue these isomorphisms to an isomorphism
between $X$ and $\underline{\text{Proj}}_S(\mathcal{A})$ as
desired. Details omitted.

\medskip\noindent
Here is the second argument.
Lemma \ref{lemma-glue-relative-proj-twists}
shows that there exists a morphism of graded algebras
$$
\psi : \pi^*\mathcal{A}
\longrightarrow
\bigoplus\nolimits_{n \geq 0} \mathcal{O}_Y(n)
$$
over $Y$ which on sections over affine opens of $S$ agrees with
(\ref{equation-global-sections}). Hence for every $y \in Y$
there exists an open neighbourhood $V \subset Y$ of $y$
and an integer $d \geq 1$ such that for $d | n$ the sheaf
$\mathcal{O}_Y(n)|_V$ is invertible and the multiplication maps
$\mathcal{O}_Y(n)|_V \otimes_{\mathcal{O}_V} \mathcal{O}_Y(m)|_V
\to \mathcal{O}_Y(n + m)|_V$ are isomorphisms. Thus
$\psi$ restricted to the sheaf $\pi^*\mathcal{A}^{(d)}|_V$
gives an element of $F_d(V)$. Since the opens $V$ cover $Y$
we see ``$\psi$'' gives rise to an element of $F(Y)$.
Hence a canonical morphism $Y \to X$ over $S$.
Because this construction is completely canonical to see
that it is an isomorphism we may work locally on $S$.
Hence we reduce to the case $S$ affine where the result is clear.
\end{proof}

\begin{definition}
\label{definition-relative-proj}
Let $S$ be a scheme. Let $\mathcal{A}$ be a quasi-coherent sheaf of
graded $\mathcal{O}_S$-algebras. The
{\it relative homogeneous spectrum of $\mathcal{A}$ over $S$},
or the {\it homogeneous spectrum of $\mathcal{A}$ over $S$}, or the
{\it relative Proj of $\mathcal{A}$ over $S$} is the scheme
constructed in Lemma \ref{lemma-glue-relative-proj} which represents the
functor $F$ (\ref{equation-proj}), see
Lemma \ref{lemma-glueing-gives-functor-proj}.
We denote it $\pi : \underline{\text{Proj}}_S(\mathcal{A}) \to S$.
\end{definition}

\noindent
The relative Proj comes equipped with a quasi-coherent
sheaf of $\mathbf{Z}$-graded algebras
$\bigoplus_{n \in \mathbf{Z}}
\mathcal{O}_{\underline{\text{Proj}}_S(\mathcal{A})}(n)$
(the twists of the structure sheaf) and
a ``universal'' homomorphism of graded algebras
$$
\psi_{univ} :
\mathcal{A}
\longrightarrow
\pi_*\left(
\bigoplus\nolimits_{n \geq 0}
\mathcal{O}_{\underline{\text{Proj}}_S(\mathcal{A})}(n)
\right)
$$
see Lemma \ref{lemma-glue-relative-proj-twists}. We may also think of this
as a homomorphism
$$
\psi_{univ} :
\pi^*\mathcal{A}
\longrightarrow
\bigoplus\nolimits_{n \geq 0}
\mathcal{O}_{\underline{\text{Proj}}_S(\mathcal{A})}(n)
$$
if we like. The following lemma is a formulation of the
universality of this object.

\begin{lemma}
\label{lemma-tie-up-psi}
In Situation \ref{situation-relative-proj}.
Let $(f : T \to S, d, \mathcal{L}, \psi)$
be a quadruple. Let
$r_{d, \mathcal{L}, \psi} : T \to \underline{\text{Proj}}_S(\mathcal{A})$
be the associated $S$-morphism.
There exists an isomorphism
of $\mathbf{Z}$-graded $\mathcal{O}_T$-algebras
$$
\theta :
r_{d, \mathcal{L}, \psi}^*\left(
\bigoplus\nolimits_{n \in \mathbf{Z}}
\mathcal{O}_{\underline{\text{Proj}}_S(\mathcal{A})}(nd)
\right)
\longrightarrow
\bigoplus\nolimits_{n \in \mathbf{Z}} \mathcal{L}^{\otimes n}
$$
such that the following diagram commutes
$$
\xymatrix{
\mathcal{A}^{(d)} \ar[rr]_-{\psi}
 \ar[rd]_-{\psi_{univ}} & &
f_*\left(
\bigoplus\nolimits_{n \in \mathbf{Z}}
\mathcal{L}^{\otimes n}
\right) \\
 &
\pi_*\left(
\bigoplus\nolimits_{n \geq 0}
\mathcal{O}_{\underline{\text{Proj}}_S(\mathcal{A})}(nd)
\right) \ar[ru]_\theta
}
$$
The commutativity of this diagram uniquely determines $\theta$.
\end{lemma}

\begin{proof}
Note that the quadruple $(f : T \to S, d, \mathcal{L}, \psi)$
defines an element of $F_d(T)$. Let
$U_d \subset \underline{\text{Proj}}_S(\mathcal{A})$
be the locus
where the sheaf $\mathcal{O}_{\underline{\text{Proj}}_S(\mathcal{A})}(d)$
is invertible and generated by the image of
$\psi_{univ} : \pi^*\mathcal{A}_d \to
\mathcal{O}_{\underline{\text{Proj}}_S(\mathcal{A})}(d)$.
Recall that $U_d$ represents the functor $F_d$, see the proof
of Lemma \ref{lemma-relative-proj}. Hence the result will follow
if we can show the quadruple
$(U_d \to S, d, \mathcal{O}_{U_d}(d), \psi_{univ}|_{\mathcal{A}^{(d)}})$
is the universal family, i.e., the representing object in $F_d(U_d)$.
We may do this after restricting to an affine open of $S$ because
(a) the formation of the functors $F_d$ commutes with base change
(see Lemma \ref{lemma-proj-base-change}), and (b) the pair
$(\bigoplus_{n \in \mathbf{Z}}
\mathcal{O}_{\underline{\text{Proj}}_S(\mathcal{A})}(n),
\psi_{univ})$
is constructed by glueing over affine opens in $S$
(see Lemma \ref{lemma-glue-relative-proj-twists}).
Hence we may assume that $S$ is affine. In this case the functor
of quadruples $F_d$ and the functor of triples $F_d$ agree
(see proof of Lemma \ref{lemma-relative-proj-affine}) and moreover
Lemma \ref{lemma-proj-functor-strict}
shows that $(d, \mathcal{O}_{U_d}(d), \psi^d)$
is the universal triple over $U_d$.
Going backwards through the identifications in the proof of
Lemma \ref{lemma-relative-proj-affine} shows that
$(U_d \to S, d, \mathcal{O}_{U_d}(d), \psi_{univ}|_{\mathcal{A}^{(d)}})$
is the universal quadruple as desired.
\end{proof}

\begin{lemma}
\label{lemma-relative-proj-separated}
Let $S$ be a scheme and $\mathcal{A}$ be a quasi-coherent sheaf
of graded $\mathcal{O}_S$-algebras. The morphism
$\pi : \underline{\text{Proj}}_S(\mathcal{A}) \to S$
is separated.
\end{lemma}

\begin{proof}
To prove a morphism is separated we may work locally on the base,
see Schemes, Section \ref{schemes-section-separation-axioms}.
By construction $\underline{\text{Proj}}_S(\mathcal{A})$ is
over any affine $U \subset S$ isomorphic to
$\text{Proj}(A)$ with $A = \mathcal{A}(U)$. By
Lemma \ref{lemma-proj-separated} we see that $\text{Proj}(A)$ is separated.
Hence $\text{Proj}(A) \to U$ is separated (see
Schemes, Lemma \ref{schemes-lemma-compose-after-separated}) as desired.
\end{proof}

\begin{lemma}
\label{lemma-relative-proj-base-change}
Let $S$ be a scheme and $\mathcal{A}$ be a quasi-coherent sheaf
of graded $\mathcal{O}_S$-algebras. Let $g : S' \to S$ be any morphism
of schemes. Then there is a canonical isomorphism
$$
r :
\underline{\text{Proj}}_{S'}(g^*\mathcal{A})
\longrightarrow
S' \times_S \underline{\text{Proj}}_S(\mathcal{A})
$$
as well as a corresponding isomorphism
$$
\theta :
r^*\text{pr}_2^*\left(\bigoplus\nolimits_{d \in \mathbf{Z}}
\mathcal{O}_{\underline{\text{Proj}}_S(\mathcal{A})}(d)\right)
\longrightarrow
\bigoplus\nolimits_{d \in \mathbf{Z}}
\mathcal{O}_{\underline{\text{Proj}}_{S'}(g^*\mathcal{A})}(d)
$$
of $\mathbf{Z}$-graded
$\mathcal{O}_{\underline{\text{Proj}}_{S'}(g^*\mathcal{A})}$-algebras.
\end{lemma}

\begin{proof}
This follows from Lemma \ref{lemma-proj-base-change} and the construction
of $\underline{\text{Proj}}_S(\mathcal{A})$ in
Lemma \ref{lemma-relative-proj} as the union
of the schemes $U_d$ representing the functors $F_d$.
In terms of the construction of relative Proj via glueing
this isomorphism is given by the isomorphisms constructed
in Lemma \ref{lemma-base-change-map-proj} which provides us with
the isomorphism $\theta$. Some details omitted.
\end{proof}

\begin{lemma}
\label{lemma-apply-relative}
Let $S$ be a scheme.
Let $\mathcal{A}$ be a quasi-coherent sheaf of graded $\mathcal{O}_S$-modules
generated as an $\mathcal{A}_0$-algebra by $\mathcal{A}_1$.
In this case the scheme $X = \underline{\text{Proj}}_S(\mathcal{A})$
represents the functor $F_1$ which associates to a scheme
$f : T \to S$ over $S$ the set of pairs $(\mathcal{L}, \psi)$, where
\begin{enumerate}
\item $\mathcal{L}$ is an invertible $\mathcal{O}_T$-module, and
\item $\psi : f^*\mathcal{A} \to \bigoplus_{n \geq 0} \mathcal{L}^{\otimes n}$
is a graded $\mathcal{O}_T$-algebra homomorphism such that
$f^*\mathcal{A}_1 \to \mathcal{L}$ is surjective
\end{enumerate}
up to strict equivalence as above. Moreover, in this case all the
quasi-coherent sheaves
$\mathcal{O}_{\underline{\text{Proj}}(\mathcal{A})}(n)$
are invertible
$\mathcal{O}_{\underline{\text{Proj}}(\mathcal{A})}$-modules
and the multiplication maps induce isomorphisms
$
\mathcal{O}_{\underline{\text{Proj}}(\mathcal{A})}(n)
\otimes_{\mathcal{O}_{\underline{\text{Proj}}(\mathcal{A})}}
\mathcal{O}_{\underline{\text{Proj}}(\mathcal{A})}(m) =
\mathcal{O}_{\underline{\text{Proj}}(\mathcal{A})}(n + m)$.
\end{lemma}

\begin{proof}
Under the assumptions of the lemma the sheaves
$\mathcal{O}_{\underline{\text{Proj}}(\mathcal{A})}(n)$
are invertible and the multiplication maps isomorphisms
by Lemma \ref{lemma-relative-proj} and
Lemma \ref{lemma-apply}
over affine opens of $S$. Thus $X$ actually represents the
functor $F_1$, see proof of Lemma \ref{lemma-relative-proj}.
\end{proof}












